\section{Encoder đọc cả hai chiều}
\label{sec:ch06-hai-chieu}

\index{encoder hai chiều!lý do bài đưa ra}%
Còn một nửa của kiến trúc tôi chưa đụng tới. Mục 3.2 của bài thay encoder một
chiều bằng \tn{encoder hai chiều}{bidirectional encoder}: hai mạng hồi quy,
một đọc câu từ trái sang phải, một đọc từ phải sang trái, và annotation của vị
trí \(j\) là hai trạng thái nối lại,
\(\vect{h}_j = [\overrightarrow{\vect{h}}_j ; \overleftarrow{\vect{h}}_j]\),
tức phương trình (7). Ý tưởng gốc là của Schuster và
Paliwal~\autocite{schuster1997bidirectional}.

Lý do bài đưa ra là ngữ nghĩa: họ muốn annotation tóm tắt \enquote{not only the
preceding words, but also the following words}. Với encoder một chiều,
\(\vect{h}_j\) chỉ biết đoạn câu từ đầu tới \(j\), nên nó không thể biết từ
\(j\) sắp được dùng làm gì.

Mục này hỏi hai chuyện, và cả hai đều là chuyện đo được.

\subsection*{Câu hỏi thứ nhất: phép đảo câu nguồn còn tác dụng không}

\index{đảo câu nguồn!sau khi có attention}%
Mục~\ref{sec:ch05-dao-cau-nguon} đo được rằng đảo câu nguồn đem lại 0.1789 điểm
khớp đúng cả câu, và toàn bộ phần lợi ấy rơi vào câu dài. Lập luận của
Sutskever và cộng sự là lập luận về khoảng cách: đảo câu làm những từ nguồn đầu
tiên nằm sát những từ đích đầu tiên, nên rút ngắn quãng ngắn nhất mà tín hiệu
lỗi phải đi.

Lập luận ấy đưa ra một dự đoán về chương này, và dự đoán ấy sai được. Nếu phép
đảo có tác dụng vì nó rút ngắn một đường dài, thì trong một mô hình mà đường
dài đã được thay bằng đường một bước, nó phải hết tác dụng.

\subsection*{Câu hỏi thứ hai: chiều ngược đáng giá bao nhiêu}

\index{encoder hai chiều!phép cắt bỏ chiều ngược}%
Bỏ mạng đọc ngược đi, giữ nguyên mọi thứ khác. Vẫn một annotation cho mỗi vị
trí, vẫn tổng có trọng số, vẫn phương trình~\eqref{eq:ch06-tong-co-trong-so}.
Chỉ là annotation không còn biết gì về đoạn câu phía sau nó. Đây không phải
quay lại mô hình chương trước; nó là mô hình chương này với đúng một bộ phận bị
tháo.

\begin{measured}{encoder hai chiều, \tn{phép cắt bỏ}{ablation}, và một dòng đối chứng}
\begin{minted}{text}
encoder        d_h  source    params  exact   short   long
bidirectional  32   raw       40805   0.9967  1.0000  0.9934
bidirectional  32   reversed  40805   1.0000  1.0000  1.0000
forward only   32   raw       27365   0.9500  0.9932  0.9079
forward only   32   reversed  27365   0.8667  0.9932  0.7434
forward only   42   raw       41495   1.0000  1.0000  1.0000
\end{minted}

Cùng thiết lập với các bảng trước, giải mã tham lam, một seed mỗi dòng.

Đo bằng \code{experiments/ch06\_encoder.py} tại tag \code{ch06}.
\end{measured}

\index{đảo câu nguồn!hết tác dụng}%
Câu hỏi thứ nhất trả lời xong ở hai dòng đầu. Với encoder hai chiều, đảo câu
nguồn đem lại 0.0033, tức một câu trong ba trăm. Chương trước đo được 0.1789
trung bình trên ba seed cho cùng phép biến đổi ấy trên cùng corpus ấy.

Dự đoán đúng, và nó đúng theo cách đáng tin nhất: nó là dự đoán rút ra từ lời
giải thích của \emph{bài trước}, kiểm bằng mô hình của \emph{bài này}. Nếu phép
đảo vẫn đem lại mười mấy điểm ở đây thì lời giải thích theo khoảng cách của
Sutskever và cộng sự đã sai, dù con số của họ vẫn đúng.

\subsection*{Dòng cuối là dòng đổi câu trả lời}

\index{encoder hai chiều!chỗ tôi suýt kết luận sai}%
Câu hỏi thứ hai thì tôi suýt trả lời sai, và tôi kể lại vì cái sai ấy rẻ tiền
và ai cũng mắc được.

Đọc bốn dòng đầu thì rõ ràng: hai chiều được 0.9967, một chiều được 0.9500, và
chỗ mất nằm ở câu dài, 0.9934 xuống 0.9079. Kết luận có sẵn: chiều ngược đáng
giá gần năm điểm. Tôi đã viết câu đó ra rồi.

Rồi tôi đếm tham số. Mô hình hai chiều có 40805, mô hình một chiều có 27365.
Bỏ một mạng hồi quy đi thì mất một phần ba số tham số, nên phép so ấy trộn lẫn
\enquote{đọc được cả hai phía} với \enquote{to hơn}. Đúng cái bẫy mà
mục~\ref{sec:ch06-be-rong} vừa phải tránh một lần.

Dòng cuối là phép đối chứng. Encoder một chiều ở \(d_h = 42\) mang 41495 tham
số, hơn mô hình hai chiều 1.7 phần trăm, và nó được \textbf{1.0000}. Trên cả ba
cột.

Vậy trên corpus này, chiều ngược không đáng đồng nào. Phần hơn trong bốn dòng
đầu là phần tham số.

\subsection*{Chuyện này bác corpus, không bác bài báo}

\index{corpus đồ chơi!không có chỗ cho ngữ cảnh phải}%
Đây là chỗ dễ đọc sai nhất trong cả chương, nên tôi nói thẳng: dòng cuối
\emph{không} phải bằng chứng rằng mục 3.2 của bài sai.

Lý do bài đưa ra cho chiều ngược là ngữ nghĩa: cần biết các từ đứng sau để hiểu
từ đang xét. Mà văn phạm sinh ra corpus này không có một chỗ nhập nhằng nào.
Mục~\ref{sec:ch05-corpus} đã liệt kê điều đó ngay từ chương trước, trong danh
sách những thứ corpus không mô hình hóa: không hình thái, không hợp giống số,
không nhập nhằng, không từ hiếm. Mỗi từ tiếng Anh ở đây ánh xạ sang đúng một
thứ bất kể ngữ cảnh phải là gì.

Một corpus không có nhập nhằng thì không thể cho thấy công dụng của một cơ chế
sinh ra để gỡ nhập nhằng. Phép đo này đo corpus chứ không đo bài báo, và cách
duy nhất để nó nói được điều gì về bài báo là chạy trên dữ liệu có tính chất mà
cơ chế nhắm tới. Bài tập tier ba có một câu hỏi dựng đúng thí nghiệm ấy.

Cái tôi rút ra được, và nó vẫn đáng giá, là một chuyện về phương pháp: một phép
cắt bỏ không kèm dòng đối chứng về tham số thì đo được rất ít. Bốn dòng đầu tự
chúng trông thuyết phục.

\subsection*{Một chỗ tôi chưa giải thích được}

\index{encoder hai chiều!đảo câu làm hại encoder một chiều}%
Dòng thứ tư của bảng làm tôi khó chịu và tôi để nó lại như một quan sát.

Với encoder một chiều, đảo câu nguồn \emph{làm hại}: 0.9500 xuống 0.8667. Chỗ
mất nằm gọn ở câu dài, 0.9079 xuống 0.7434, còn câu ngắn đứng nguyên 0.9932 ở
cả hai dòng.

Tôi có vài cách giải thích nghe hợp lý và không kiểm được cái nào trong ngân
sách chương này, nên tôi không in cái nào ra. Điều đáng nói là chiều của nó:
với encoder hai chiều phép đảo gần như vô hại, với encoder một chiều nó có hại,
nên tương tác giữa hai thứ ấy là thật chứ không phải nhiễu. Bài tập cuối chương
hỏi lại câu này.

\subsection*{Cả hai chiều dừng ở đâu, và vì sao chuyện đó dễ sai}

\index{encoder hai chiều!chỗ dừng ở token đệm}%
Một chi tiết cài đặt, vì nó thuộc loại hỏng mà không có gì báo.

Trong một batch, câu ngắn được đệm cho bằng câu dài nhất. Chiều xuôi thì
chương~\ref{ch:encoder-decoder} đã xử lý: đóng băng trạng thái ở mọi vị trí
đệm, nên câu ngắn kết thúc ở trạng thái tại từ cuối của chính nó. Chiều ngược
gặp đúng bài toán ấy ở đầu kia. Chạy nó từ cột cuối cùng của batch thì mọi câu
ngắn bắt đầu bằng một loạt bước trên token đệm.

\begin{minted}{python}
ahead = self._scan(self.ahead, embedded, live, range(steps))
if self.back is None:
    return torch.stack(ahead), mask

backwards = reversed(range(steps))
back = self._scan(self.back, embedded, live, backwards)
stacked = [torch.stack(ahead), torch.stack(back)]
return torch.cat(stacked, dim=-1), mask
\end{minted}

Cùng một hàm \code{\_scan} chạy cho cả hai chiều, khác nhau đúng thứ tự duyệt,
và phép đóng băng nằm bên trong nó:

\begin{minted}{python}
for t in order:
    c_next, h_next = layer.step_projected(projected[t], state)
    state = (
        live[t] * c_next + (1.0 - live[t]) * state[0],
        live[t] * h_next + (1.0 - live[t]) * state[1],
    )
    out[t] = state[1]
\end{minted}

Nhờ phép đóng băng ấy mà chiều ngược tự bắt đầu ở từ cuối của mỗi câu: ở các
vị trí đệm nó giữ nguyên trạng thái khởi tạo, nên bước cập nhật thật đầu tiên
của nó rơi đúng vào token thật cuối cùng. Repo có test khẳng định annotation
của một câu không đổi khi đem câu ấy xếp chung batch với câu dài hơn.
